\chapter{引言}

\section{研究背景}

学位论文应说明选题意义、国内外研究现状、拟解决的问题、采用的主要理论与方法以及论文结构。本文档中的文字仅用于演示模板结构，正式写作时应替换为论文实际内容。

参考文献引用采用顺序编码制，例如国标文献著录规则可写作\cite{gbt7714-2015}，与模板相关的说明性文献也可按引用顺序自动编号\cite{lamport1994latex}。本模板还给出中文期刊论文\cite{zhang2024article}、中文会议论文\cite{li2023conference}、硕士学位论文\cite{wang2022master}、博士学位论文\cite{chen2021phd}、英文期刊论文\cite{smith2024englisharticle}和英文会议论文\cite{johnson2023englishconference}等常见文献类型的著录示例。

\section{研究内容}

本模板示例覆盖以下内容：正文段落、标题层级、图题、表题、公式编号和参考文献。学校规范不建议无特殊情况使用三级及以上节标题，因此模板只给出到二级节标题的常用示例。

\subsection{模板结构}

模板入口文件为 \verb|main.tex|，版式定义集中在 \verb|ussthesis.cls|，正文内容位于 \verb|chapters| 目录。用户通常只需要修改入口文件中的论文信息和各章节内容。

\section{论文组织}

第一章为引言，第二章演示正文元素，第三章为结论。参考文献之后可以根据需要加入附录、在读期间成果和致谢。

偶数页测试。

偶数页测试。

偶数页测试。

偶数页测试。

偶数页测试。

偶数页测试。

偶数页测试。

偶数页测试。

偶数页测试。

偶数页测试。

偶数页测试。

偶数页测试。

偶数页测试。

偶数页测试。